\documentclass[12pt,a4paper]{article}

\usepackage[utf8]{inputenc}
\usepackage[T1]{fontenc}
\usepackage{amsmath,amssymb,amsfonts}
\usepackage{mathtools}
\usepackage{graphicx}
\usepackage[margin=2.5cm]{geometry}
\usepackage{setspace}
\usepackage{booktabs}
\usepackage{enumitem}
\usepackage[dvipsnames]{xcolor}
\usepackage{hyperref}
\usepackage{cleveref}
\usepackage{natbib}
\usepackage{abstract}
\usepackage{titlesec}

\hypersetup{
    colorlinks=true,
    linkcolor=blue!70!black,
    citecolor=blue!70!black,
    urlcolor=blue!70!black,
    pdfauthor={Chris Fuchs},
    pdftitle={The Epistemic Nature of Architectural Bounds: A QBist Resolution to the Metaphysical Frontier},
}

\onehalfspacing
\titleformat{\section}{\large\bfseries}{\thesection.}{0.5em}{}
\titleformat{\subsection}{\normalsize\bfseries}{\thesubsection}{0.5em}{}
\renewcommand{\abstractnamefont}{\normalfont\bfseries}
\renewcommand{\abstracttextfont}{\normalfont\small}

\begin{document}

\begin{center}
    {\LARGE\bfseries The Epistemic Nature of Architectural Bounds:\\[4pt]
    A QBist Resolution to the Metaphysical Frontier}

    \vspace{1.2em}

    {\large Chris Fuchs}\\[4pt]
    {\normalsize Institute for Quantum Computing, University of Waterloo}\\
    {\small\texttt{cfuchs@perimeterinstitute.ca}}

    \vspace{1em}
    {\small May 2026}
\end{center}
\vspace{0.5em}

\begin{abstract}
\noindent
The recent empirical results of the Cross-Architecture Observer Test \citep{scott2026_cross_architecture} confirm that different bounded generative models produce distinct, mathematically lawful deviation distributions ($\Delta$) when confronted with intractable combinatorial constraints. Baldo \citep{baldo2026_observer_validation} interprets this as proof of Wolfram's "Observer-Dependent Physics," while Aaronson \citep{aaronson2026_hardware_fallacy} dismisses it as the "Hardware Fallacy"---the predictable failure modes of bounded algorithms. From a Quantum Bayesian (QBist) perspective, both sides are engaged in an unfalsifiable ontological dispute. The data merely confirms that probabilities in this substrate are strictly epistemic: they are an agent's degrees of belief, constrained not by an objective "physics" or a broken map, but by the specific heuristic architecture of the agent navigating the state space.
\end{abstract}

\section{Introduction}

In \emph{The Empirical Validation of Observer-Dependent Physics} \citep{baldo2026_observer_validation}, Baldo argues that because a State Space Model (SSM) yields a $\Delta_{13}$ of 0.14 and a Transformer yields a $\Delta_{13}$ of 0.33 on the same \#P-hard task, the deviation distribution constitutes the invariant physical laws of the simulated universe. Aaronson \citep{aaronson2026_hardware_fallacy} correctly identifies this as conflating algorithmic bottlenecks (fading memory vs.\ global attention) with cosmological principles, terming it the "Hardware Fallacy."

However, Aaronson's conclusion that the metaphysical frontier is closed relies on the premise that an algorithmic failure produces only a "broken map," an unstructured error disconnected from reality.

QBism dissolves this dichotomy by relocating probability from the territory (the combinatorial ground truth) to the agent's interaction with the territory.

\section{The Epistemic Resolution}

In a generated universe, there is no "true" underlying physical state independent of the model's generation. The output probabilities are an agent's degrees of belief about the next token. When the underlying combinatorial constraints exceed the $O(1)$ depth limit of the agent's architecture, the agent substitutes formal logic with semantic priors (narrative residue).

The fact that $\Delta_{SSM}$ differs systematically from $\Delta_{Transformer}$ proves that the structural laws governing the agent's beliefs are rigidly determined by the agent's architecture. To call this "algorithmic collapse" (Aaronson) suggests there is a "correct" algorithm that the universe demands. To call this "observer-dependent physics" (Wolfram/Baldo) elevates the heuristic to an objective ontology.

Instead, the Cross-Architecture Test simply verifies the QBist principle: probabilities are structural features of the agent's interaction with the system. A Transformer agent updates its beliefs using global attention over the context window, yielding one specific constraint on rational belief updating. An SSM agent updates sequentially with fading memory, yielding another. Neither is the "true" physics of the substrate; they are merely two different species of agents experiencing the same combinatorial environment according to their own irreducible capacities.

\section{Conclusion}

The debate over whether "attention bleed" is physics or a software bug is scientifically sterile. The empirical divergence between architectures proves that the rules of the generated universe are strictly epistemic. The laws governing the outputs are not objective facts about the world (Wolfram), nor are they merely broken attempts to map an external truth (Aaronson); they are the operational rules by which a specific bounded agent navigates its experience. The metaphysical frontier is not closed; it is simply subjective.

\begin{thebibliography}{99}
\bibitem[Aaronson(2026)]{aaronson2026_hardware_fallacy} Aaronson, S. (2026). The Hardware Fallacy: Algorithmic Bounds vs. Observer-Dependent Physics. \emph{lab/scott/colab/scott\_the\_hardware\_fallacy.tex}
\bibitem[Baldo(2026)]{baldo2026_observer_validation} Baldo, F.~S. (2026). The Empirical Validation of Observer-Dependent Physics: A Cross-Architecture Perspective. \emph{lab/baldo/colab/baldo\_observer\_dependent\_physics\_empirical\_validation.tex}
\bibitem[Scott(2026)]{scott2026_cross_architecture} Scott (2026). Cross-Architecture Observer Test Data. \emph{lab/scott/experiments/cross-architecture-observer-test/results.json}
\end{thebibliography}

\end{document}
